This model utilizes the \textbf{DistilBERT}-based classification architecture provided by Hugging Face's \texttt{AutoModelForSequenceClassification}. Most importantly, it utilizes a base pretrained version of the model known as \textbf{DistilBERT-base-uncased}. This transformer model was chosen in favour of the complete base version of BERT due to its \textbf{decreased size, computational complexity, and training cost}. Despite this, the model has been proven to retain a \textbf{strong contextual understanding of textual features} during usage. Specifically, the model uses a \textbf{6-layer encoder architecture} compared to the \textbf{12-layer architecture} of the base BERT model. After model importing, it was configured for the \textbf{binary classification task} (song genre classification, Pop and Metal classes) using \textbf{raw lyric textual features} and utilized a standard label mapping of \texttt{p: 0} and \texttt{m: 1} (for easier mapping during downstream training). Initial preprocessing of features (raw lyric text) only involved \textbf{normalizing whitespace} and \textbf{removing null/empty textual records} to retain the original lexical content of the dataset. Afterward, tokenization was conducted using the pretrained \textbf{DistilBERT tokenizer}, which relies on \textbf{WordPiece subword tokenization} (decomposes words into smaller words). WordPiece subword tokenization decomposes a single word into \textbf{multiple smaller tokens} for better size and efficiency. Before passing into the tokenizer, all textual data sequences were \textbf{truncated to a maximum length of 256 tokens}. A token length of \textbf{256} was chosen to allow for \textbf{training and computational efficiency} without largely undermining the \textbf{contextual analysis of lyrics} within the model.

For model fine-tuning, the parameters involved \textbf{2 epochs}, \textbf{input batch size: 16}, \textbf{learning rate: 0.00002}, \textbf{weight decay: 0.01}, and \textbf{warmup steps: 300}. Moreover, during training, the \textbf{Adam optimizer} along with \textbf{cross-entropy loss} (loss function methodology) was used, followed by validation (using the validation set) at the end of each epoch. During training, the \textbf{best-performing checkpoint} was selected using the \textbf{Macro F1 score}. As a result, the final model achieved an \textbf{Accuracy and Macro F1 score of 86.46\%} (\textbf{Accuracy: 0.8646, Macro-F1: 0.8646}) using the test set. Based on the prior baseline performances (baseline performances in Table 1), the \textbf{DistilBERT model outperforms the majority baseline} and \textbf{slightly improves upon the LR + TF-IDF baseline}. This conveys that even though \textbf{class-based lexical classification of lyrics is already strong} (from the LR-based baseline), the introduction of \textbf{contextual semantic text representations} from the DistilBERT model may provide \textbf{slight performance enhancements}.